% Introduction (draft). Numbers from paper/tables (scripts/build_paper.py) and results/stats/a3_stall.csv.
\section{Introduction}
\label{sec:intro}

Mobile and Internet-of-Things devices run applications whose computation exceeds what the device can deliver within
its time and energy budget. Offloading parts of an application to edge servers close to the device, or further to a
cloud data centre, shortens execution but costs transmission time, energy and money, so where each part runs is a
multi-objective decision. Two features of real deployments make it hard. First, applications are rarely a single
task: a video analytics pipeline, a scientific workflow or a multi-stage inference service is a directed acyclic graph
(DAG) of tasks that exchange data, and the completion time is set by the longest chain of dependent computation and
transfer, not by any single task. Second, the quantities the decision depends on are uncertain at planning time:
wireless bandwidth fluctuates, servers are shared, and task workloads vary with their input. When only bounds on
these quantities are known, each objective becomes an interval, and the decision problem an interval multi-objective
optimisation problem.

Interval multi-objective evolutionary algorithms handle such problems by comparing solutions through interval
dominance and by rewarding solutions whose objective intervals are narrow \citep{limbourg2005ipmoea,jin2023imoeadawn,cui2024imeaied}.
The most recent of them for offloading, IMOEA-DQN \citep{jin2026imoeadqn}, selects each new population in three tiers:
a share $\alpha$ of the places by convergence (an interval dominance degree), a share $\beta$ by diversity (an interval
crowding distance) and a share $\gamma$ by uncertainty, and lets a deep Q-network choose $(\alpha, \beta, \gamma)$ in every
generation. It was formulated for independent tasks in a cloud--edge--end system with interval resources.

This paper extends that line of work to DAG applications and, in doing so, finds that two of its parts do not
carry over. The interval dominance degree loses its selection pressure: once most of the merged population is mutually
non-dominated, more pool members share the best value than there are convergence places, and the tier is filled by the
tie-break, at random. On our instances this happens for 86--100\% of the convergence places from generation 100 on, and
at the end of a run 117--121 of the 210 pool members share the best value for 88 places. On the 1000-task workflows the
backbone of IMOEA-DQN, run with the same fixed action as our method, ends with an IGD about six times higher. And the
learned controller of $(\alpha, \beta, \gamma)$ does not beat a well-chosen fixed action, because the per-generation reward it
learns from is almost pure noise.

We propose \ours{}, an interval multi-objective evolutionary algorithm for offloading DAG workflows in uncertain
cloud--edge--end environments. Its contributions are:
\begin{enumerate}
\item An interval critical-path model of DAG offloading over devices, edge servers and the cloud, with makespan, device
  energy and cost as interval objectives, task placement and task order as decision variables, and data transfer along
  every dependency. For single-task applications it reduces exactly to the model of \citet{jin2026imoeadqn}
  (Section~\ref{sec:model}).
\item A DAG-valid two-string encoding with precedence-preserving crossover and mutation, so that every offspring is
  feasible without repair (Section~\ref{sec:variation}).
\item A diagnosis of why the interval dominance degree stalls as a convergence tier, and a fix: an interval version of
  the best-position Tchebycheff rank of MSOPS \citep{hughes2003msops}, with ties broken by crowding distance, used as the
  convergence tier inside the three-tier selection, together with mating among neighbours in objective space
  (Sections~\ref{sec:stall} and~\ref{sec:rankc}). The rank itself is not new; its use inside the dominance-, diversity-
  and uncertainty-aware selection, the diagnosis and the evidence are.
\item An analysis of learned control of the selection weights: three learner designs with and without a
  graph-attention encoder of the DAG, and a reward audit that explains why none of them beats a fixed action
  (Section~\ref{sec:learned}).
\item An evaluation on 57 instances (Jin et al.'s instance sizes and five scientific workflow families with up to 1000
  tasks), plus 15 instances with 2000 tasks for scalability, against four interval evolutionary algorithms, NSGA-II,
  NSGA-III, IMOEA-DQN and list-scheduling heuristics, with a robustness check of the normalisation, a Monte Carlo check
  of the interval objectives against realised values, and an account of where the method is weaker
  (Sections~\ref{sec:experiments} and~\ref{sec:discussion}).
\end{enumerate}

On Jin et al.'s instance sizes \ours{} has the best IGD and hypervolume on all 12 instances, with 62\% lower mean IGD
than IMOEA/D. On the 45 workflow instances it is better than IMOEA/D on 43 instances in IGD and on 42 in hypervolume,
at about the same running time. At 1000 and 2000 tasks it keeps the hypervolume lead, while IMOEA/D covers the cheapest
schedules better and IGD is level.

Section~\ref{sec:related} reviews related work, Section~\ref{sec:model} states the model and
Section~\ref{sec:method} the method. Section~\ref{sec:experiments} reports the experiments,
Section~\ref{sec:discussion} discusses limitations, and Section~\ref{sec:conclusion} concludes.
